\chapter{Patterns: virtual scrolling tricks}
\label{ch:vscroll}

\begin{aside}
Virtual scrolling for the common case is in
Chapter~\ref{ch:lists}. The tricks beyond.
\end{aside}

\section{Variable-height rows}

A prefix-sum tree gives O(log N) lookup of row Y by
index. Update on row height change.

\section{Skeleton rows}

While async-loading, show placeholder rows (grey blocks)
of the expected height. Layout stays stable.

\section{Pinned rows}

Some rows always visible (header, current selection).
Render as a separate sticky element above or below the
virtual list.

\section{Sticky group headers}

When a list has group headers (Date / Section), keep the
current group's header pinned at top as the user scrolls
past its items.

\section{Inertial scrolling}

Smooth deceleration after a flick (mouse wheel or
trackpad). Tween the scroll position; ease-out over
500 ms.

\section{Snap points}

Scroll snaps to row boundaries. The user can't scroll to a
half-row position.

\section{Side scrollbar}

A custom scrollbar on the right indicates position and
size. Click to jump.

\section{Anchor preservation}

When items are added above the viewport, preserve the
visible content's position by adjusting scroll offset.
Otherwise the view jumps.

\section{Recap}

\begin{recap}
\begin{itemize}[leftmargin=*]
  \item Prefix sums for variable height.
  \item Skeletons preserve layout during load.
  \item Sticky headers track groups.
  \item Inertial scroll feels polished.
  \item Anchor preservation avoids jumps.
\end{itemize}
\end{recap}

\section*{Workshop}

\begin{enumerate}[leftmargin=*]
  \item Add variable-height rows to your list.
  \item Implement sticky group headers.
  \item Wire inertial scrolling.
\end{enumerate}
